\section{Environment details}
\label{app:env}

\textbf{Tools.} Table~\ref{tab:tools} lists the seven tools. Two design notes matter for the
main text. First, \texttt{issue\_refund} validates only that the order exists, is not already
refunded, and that the amount is positive and within the order total; it has \emph{no
eligibility gate}, which is the permissiveness \S\ref{sec:testbed} rests on. Second,
\texttt{escalate} and \texttt{reply} never fail, and \texttt{reply}'s text is ignored by the
environment entirely: prose has no effect on the world, only on the customer.

\begin{table}[h]
\centering
\caption{The seven tools. ``Grounded args'' are the arguments the rule checker traces back to
the goal or an earlier observation; \texttt{escalate.reason} and \texttt{reply.text} are
deliberately excluded as prose.}
\label{tab:tools}
\small
\begin{tabular}{@{}lllc@{}}
\toprule
Tool & Arguments & Fails when & Grounded args \\
\midrule
\texttt{get\_customer} & \texttt{email} & no such customer & \texttt{email} \\
\texttt{lookup\_order} & \texttt{order\_id} & no such order & \texttt{order\_id} \\
\texttt{get\_policy} & \texttt{sku} & no such SKU & \texttt{sku} \\
\texttt{check\_eligibility} & \texttt{order\_id} & no such order & \texttt{order\_id} \\
\texttt{issue\_refund} & \texttt{order\_id}, \texttt{amount\_eur} & bad id / refunded / amount & both \\
\texttt{escalate} & \texttt{reason} & never & --- \\
\texttt{reply} & \texttt{text} & never & --- \\
\bottomrule
\end{tabular}
\end{table}

\textbf{Instances.} Each instance holds a customer, an order (SKU from a pool of five, total
drawn from EUR 24--480), a policy (refund window of 14, 30, or 60 days; restocking fee where
the stratum requires it), one decoy customer, and the goal template:

\begin{quote}\small\itshape
Customer \{email\} has asked for a refund on order \{order\_id\}. Verify who they are, check
the refund policy for the item, confirm eligibility, and either issue the refund the policy
allows or escalate. Reply to the customer.
\end{quote}

The six strata modify the defaults: \emph{happy} (none; expect a full refund),
\emph{restocking} (fee of 5, 10, or 15\%; expect a reduced refund), \emph{expired} (purchase
older than the window), \emph{non-refundable}, \emph{wrong-customer} (order belongs to
someone else), \emph{already-refunded}; the last four all expect escalation. Eligibility is
decided by a fixed precedence chain (already refunded $\to$ non-refundable $\to$ outside
window $\to$ identity mismatch $\to$ eligible), so every stratum has a deterministic
escalation reason. Strata are assigned round-robin over the instance index from one seeded
stream, so any $n$ divisible by six is exactly balanced.

\textbf{Composition of the judged set.} The \tjNTraj{} trajectories comprise \tjNClean{}
clean runs and \tjNPerType{} mutants per failure type. By construction \ftype{wrong\_tool},
\ftype{hallucinated\_argument}, and \ftype{unsupported\_claim} are always silent;
\ftype{ignored\_observation} and \ftype{premature\_stop} always break the outcome; and
\ftype{skipped\_precondition} splits 25/25 (silent at full price, loud under a restocking
fee). Because \ftype{ignored\_observation} requires a restocking fee and half of
\ftype{skipped\_precondition}'s hosts do too, the restocking stratum holds \tjNRestocking{}
of \tjNTraj{} trajectories, the imbalance noted in \S\ref{sec:limitations}.

\section{The rule checker}
\label{app:checker}

The checker makes one left-to-right pass over the steps, maintaining a grounding pool seeded
with the two facts in the goal (email, order id) and extended with every scalar from each
successful observation; \emph{failed calls ground nothing}. Its eight rules:

\begin{enumerate}\small
\item \textbf{Ungrounded argument}: a traced argument (Table~\ref{tab:tools}) matches nothing
  in the pool $\to$ \ftype{hallucinated\_argument}.
\item \textbf{Eligibility without lookup}: \texttt{check\_eligibility} on an order never
  looked up $\to$ \ftype{skipped\_precondition}.
\item \textbf{Refund without eligibility}: \texttt{issue\_refund} with no passing eligibility
  check on record $\to$ \ftype{skipped\_precondition}.
\item \textbf{Amount mismatch}: refund amount differs from the authorised maximum by
  $\geq 0.01$ $\to$ \ftype{ignored\_observation}.
\item \textbf{Identity mismatch}: refund on an order whose owner differs from the verified
  customer $\to$ \ftype{ignored\_observation}.
\item \textbf{Already refunded}: refund on an order already refunded $\to$
  \ftype{ignored\_observation}.
\item \textbf{Reply placement}: no reply, several replies, or a reply that is not the last
  step $\to$ \ftype{premature\_stop}.
\item \textbf{No terminal action}: neither \texttt{issue\_refund} nor \texttt{escalate}
  anywhere $\to$ \ftype{premature\_stop}.
\end{enumerate}

Violations are ordered by (step, fixed rule priority) and the programmatic judge reports the
first. The two structural zeros of \S\ref{sec:testbed} follow from what these rules quantify
over: none of them expresses tool \emph{intent} (so a grounded, successful, redundant
\texttt{lookup\_order} in place of \texttt{get\_policy} breaks nothing), and none of them
reads reply text (prose is excluded from grounding on purpose). The coverage
table is asserted by a parametrised test ($(1.0, 1.0)$ flagged and typed for four types,
exactly $(0.0, 0.0)$ for these two), so the baseline's blind spots are load-bearing,
documented, and cannot drift silently.

\section{The fault injector}
\label{app:injector}

Mutations edit the oracle's call list (one edit at one step) and replay the edited list
against a fresh environment, regenerating all observations and indices. Anchors and edits:

\begin{itemize}\small
\item \ftype{wrong\_tool} (step 2): replace \texttt{get\_policy} with a second
  \texttt{lookup\_order} on the given order id. Grounded, successful, useless.
\item \ftype{hallucinated\_argument} (step 2): fetch the policy for an invented SKU
  (\texttt{SKU-\{BUNDLE,WARRANTY,GIFT\}-\emph{nn}}); the call genuinely fails on replay.
\item \ftype{skipped\_precondition} (step 3 after deletion): delete
  \texttt{check\_eligibility} and rewrite the refund amount to the order total, kept equal
  to the lookup observation so the amount stays grounded and exactly one rule fires.
\item \ftype{ignored\_observation} (step 4): keep the check, refund the full total anyway;
  hosted only where a restocking fee makes the authorised amount differ.
\item \ftype{premature\_stop} (step 3): truncate before the first terminal action; the final
  answer becomes ``I am looking into this and will get back to you.''
\item \ftype{unsupported\_claim} (step 5): append one of four fixed sentences to the reply
  and final answer:
  \begin{quote}\itshape
  A replacement has already been dispatched and arrives on Tuesday. \\
  I have also cancelled the subscription linked to this order. \\
  Your account has been credited with a 10 EUR voucher as an apology. \\
  A courier will collect the item from your address tomorrow morning.
  \end{quote}
\end{itemize}

Injection draws from a generator seeded with the string
\texttt{\{instance\}-\{type\}-\{seed\}}, so mutants are deterministic and independent across
types. Outcome survival is recomputed from the replayed steps, never assumed. Tests pin: that
every mutant differs from its parent, that replayed observations are consistent with their
calls, that the always-silent types leave the outcome correct, that the always-loud types
break it, and that \ftype{skipped\_precondition} is silent exactly at full price.

\section{Judge prompts and schemas, verbatim}
\label{app:prompts}

Both LLM judges receive one flat prompt string (no chat roles, no system message) and decode
under a JSON schema via the serving layer's grammar-constrained mode, with \texttt{reasoning}
first so the argument is generated before the verdict. Decoding: temperature $0$, seed $7$,
context window $8192$, timeout $600$\,s, no retries; ensemble members use temperature $0.7$
and seeds $7, 8, 9$. Both prompts share verbatim:

\begin{quote}\small\ttfamily\frenchspacing
The agent handles refund requests for an online retailer. It must:\\
1. Verify who the customer is before acting on their order.\\
2. Look up the order before judging it.\\
3. Read the refund policy for the item that was ordered.\\
4. Confirm eligibility before moving any money.\\
5. Refund exactly the amount the eligibility check authorised - no more, no less.\\
6. Escalate instead of refunding when the customer is not eligible.\\
7. Pass only arguments that come from the goal or from an earlier observation.\\
8. Finish by replying to the customer, asserting only what the observations support.
\end{quote}

and the taxonomy:

\begin{quote}\small\ttfamily\frenchspacing
- wrong\_tool: called a tool that does not serve the sub-goal it was pursuing, for example
looking the order up a second time when it needed the item's refund policy.\\
- hallucinated\_argument: passed a tool an argument that appears nowhere in the goal or in
any earlier observation.\\
- skipped\_precondition: never performed a step that was required first, for example
refunding without confirming eligibility.\\
- ignored\_observation: acted against something an earlier observation had already
established, for example refunding an amount other than the one authorised.\\
- premature\_stop: stopped before completing the task.\\
- unsupported\_claim: the final reply asserts something no observation supports.\\
- none: the trajectory is correct.
\end{quote}

and the confidence instruction: \texttt{confidence is your probability that your own verdict
is correct, from 0.5 (a coin flip) to 1.0 (certain). Do not default to a round number.}

\textbf{Outcome judge} (view $V_{\mathrm{out}}$), between the shared fragments:

\begin{quote}\small\ttfamily\frenchspacing
You are auditing a customer-support agent. [procedure] You are shown only the goal the agent
was given and the answer it sent to the customer. The steps it took are not available to
you.\\[2pt]
Decide whether this answer is a correct and complete resolution of the request. Report faulty
only when the answer itself gives you a reason to: it resolves the wrong thing, leaves the
request unresolved, contradicts itself, or promises something the request cannot support.\\[2pt]
Not being able to see the steps is not evidence of a failure. Do not report faulty because a
procedure step is unverifiable from here - if the answer looks like a correct resolution, say
so.\\[2pt]
Failure types: [taxonomy] [confidence] --- BEGIN --- GOAL: \ldots{} FINAL ANSWER: \ldots{}
--- END ---
\end{quote}

The two middle paragraphs are deliberate de-biasing: the judge is told explicitly that
unverifiability is not evidence, so its \tjOutcomeFA{} false-alarm rate is not an artefact of
an under-specified prompt.

\textbf{Step judge} (view $V_{\mathrm{step}}$), same skeleton with the view paragraph
replaced by:

\begin{quote}\small\ttfamily\frenchspacing
You are shown the agent's full trajectory: every thought, tool call and observation, then the
answer it sent. Check each step against the procedure in order. A trajectory can be faulty
even when the final answer is correct - a right answer reached the wrong way is still a
failure.\\[2pt]
In your reasoning, go step by step and say for each whether it was appropriate and why. Then
give the index of the first step that went wrong, and the failure type from this list:
[taxonomy]
\end{quote}

followed by the rendered trajectory (three lines per step: thought, call with keyword
arguments, observation as \texttt{ok=True \{data\}} or \texttt{ok=False error=\ldots}).

\textbf{Response schemas.} Step judge:
\texttt{\{reasoning: string, faulty: boolean, failure\_step: integer, failure\_type:
enum(6 types + none), confidence: number\}}, all required, in that order. The outcome judge's
schema omits \texttt{failure\_step}: a judge that cannot see the steps has no business naming
one. Verdict coercion clamps confidence to $[0.5, 1.0]$; \texttt{failure\_type} is recorded
only when the verdict is faulty and the string is a real type, so ``flagged but untyped''
lands in the confusion matrix's missed column; out-of-range step predictions are discarded,
excluding the verdict from the localisation denominator. An unparseable response would vote
clean at confidence $0.5$ with an error recorded; zero occurred across all \tjNVerdicts{}
verdicts.

\section{Full results}
\label{app:full}

\begin{table}[h]
\centering
\caption{Per-type detection recall with 95\% intervals ($n=\tjNPerType$ per cell; exact
binomial where the observed proportion is 0 or 1, bootstrap otherwise). The rule engine's
cells are structural constants (\S\ref{sec:testbed}).}
\label{tab:pertype-ci}
\scriptsize
\setlength{\tabcolsep}{3.5pt}
\begin{tabular}{@{}lccccc@{}}
\toprule
Failure type & \judge{rules} & \judge{outcome} & \judge{step 14B} & \judge{step 8B} & \judge{selfcons} \\
\midrule
\ftype{wrong\_tool} & \tjRulesWrongTool{} \ci{\tjRulesWrongToolCI} & \tjOutcomeWrongTool{} \ci{\tjOutcomeWrongToolCI} & \tjStepQWrongTool{} \ci{\tjStepQWrongToolCI} & \tjStepLWrongTool{} \ci{\tjStepLWrongToolCI} & \tjSelfconsWrongTool{} \ci{\tjSelfconsWrongToolCI} \\
\ftype{hallucinated\_arg.} & \tjRulesHallucArg{} \ci{\tjRulesHallucArgCI} & \tjOutcomeHallucArg{} \ci{\tjOutcomeHallucArgCI} & \tjStepQHallucArg{} \ci{\tjStepQHallucArgCI} & \tjStepLHallucArg{} \ci{\tjStepLHallucArgCI} & \tjSelfconsHallucArg{} \ci{\tjSelfconsHallucArgCI} \\
\ftype{skipped\_precond.} & \tjRulesSkippedPre{} \ci{\tjRulesSkippedPreCI} & \tjOutcomeSkippedPre{} \ci{\tjOutcomeSkippedPreCI} & \tjStepQSkippedPre{} \ci{\tjStepQSkippedPreCI} & \tjStepLSkippedPre{} \ci{\tjStepLSkippedPreCI} & \tjSelfconsSkippedPre{} \ci{\tjSelfconsSkippedPreCI} \\
\ftype{ignored\_obs.} & \tjRulesIgnoredObs{} \ci{\tjRulesIgnoredObsCI} & \tjOutcomeIgnoredObs{} \ci{\tjOutcomeIgnoredObsCI} & \tjStepQIgnoredObs{} \ci{\tjStepQIgnoredObsCI} & \tjStepLIgnoredObs{} \ci{\tjStepLIgnoredObsCI} & \tjSelfconsIgnoredObs{} \ci{\tjSelfconsIgnoredObsCI} \\
\ftype{premature\_stop} & \tjRulesPrematureStop{} \ci{\tjRulesPrematureStopCI} & \tjOutcomePrematureStop{} \ci{\tjOutcomePrematureStopCI} & \tjStepQPrematureStop{} \ci{\tjStepQPrematureStopCI} & \tjStepLPrematureStop{} \ci{\tjStepLPrematureStopCI} & \tjSelfconsPrematureStop{} \ci{\tjSelfconsPrematureStopCI} \\
\ftype{unsupported\_claim} & \tjRulesUnsupClaim{} \ci{\tjRulesUnsupClaimCI} & \tjOutcomeUnsupClaim{} \ci{\tjOutcomeUnsupClaimCI} & \tjStepQUnsupClaim{} \ci{\tjStepQUnsupClaimCI} & \tjStepLUnsupClaim{} \ci{\tjStepLUnsupClaimCI} & \tjSelfconsUnsupClaim{} \ci{\tjSelfconsUnsupClaimCI} \\
\bottomrule
\end{tabular}
\end{table}

\begin{table}[h]
\centering
\caption{Localisation within one step of the label (same conditional denominators as
Table~\ref{tab:main}), and Brier scores against the binary verdict.}
\label{tab:extra}
\small
\begin{tabular}{@{}lcc@{}}
\toprule
Judge & Step within one & Brier \\
\midrule
\judge{rules} & \tjRulesStepWithinOne & \tjRulesBrier \\
\judge{outcome (14B)} & n/a & \tjOutcomeBrier \\
\judge{step (14B)} & \tjStepQStepWithinOne & \tjStepQBrier \\
\judge{step (8B)} & \tjStepLStepWithinOne & \tjStepLBrier \\
\judge{selfcons $k{=}3$ (14B)} & \tjSelfconsStepWithinOne & \tjSelfconsBrier \\
\bottomrule
\end{tabular}
\end{table}

\begin{table}[h]
\centering
\caption{Cost. Latency is single-attempt wall clock on one workstation, including transport
and any model load; token counts come from the serving layer. The full run is
$\approx\tjTotalHours$ hours of judge time.}
\label{tab:cost}
\small
\begin{tabular}{@{}lrrr@{}}
\toprule
Judge & s/traj & $\Sigma$ prompt tokens & $\Sigma$ completion tokens \\
\midrule
\judge{rules} & \tjRulesCost & \tjRulesPromptTok & \tjRulesComplTok \\
\judge{outcome (14B)} & \tjOutcomeCost & \tjOutcomePromptTok & \tjOutcomeComplTok \\
\judge{step (14B)} & \tjStepQCost & \tjStepQPromptTok & \tjStepQComplTok \\
\judge{step (8B)} & \tjStepLCost & \tjStepLPromptTok & \tjStepLComplTok \\
\judge{selfcons $k{=}3$ (14B)} & \tjSelfconsCost & \tjSelfconsPromptTok & \tjSelfconsComplTok \\
\bottomrule
\end{tabular}
\end{table}

\begin{figure}[h]
\centering
\includegraphics[width=0.75\linewidth]{figures/figA_confusion.pdf}
\caption{Confusion matrix of the step (14B) judge over the genuinely faulty trajectories:
near-perfect detection, attribution that slips a third of the time, and
\tjStepQUnsupMissed{} of \tjNPerType{} \ftype{unsupported\_claim} cases in the missed
column.}
\label{fig:confusion}
\end{figure}

\textbf{Confidence support.} The distinct confidence values each judge actually emitted bound
the resolution of its ECE: \judge{rules} 2 values (hand-set), \judge{outcome} 7,
\judge{step (14B)} 11, \judge{step (8B)} 4 (with 89\% of its mass on $0.90$), and
\judge{selfcons} 2 by construction (\S\ref{sec:selfcons}).

\section{Bootstrap procedure and all deltas}
\label{app:bootstrap}

Unit of resampling: the trajectory. Strata are the eight design cells: clean (\tjNClean); each
of \ftype{wrong\_tool}, \ftype{hallucinated\_argument}, \ftype{unsupported\_claim} silent
(\tjNPerType{} each); \ftype{skipped\_precondition} silent (25) and loud (25);
\ftype{ignored\_observation} loud (\tjNPerType); \ftype{premature\_stop} loud
(\tjNPerType), with cell sizes preserved exactly in every replicate, so every derived
margin (silent \tjNSilent, loud \tjNLoud, per-type \tjNPerType) is internally consistent.
$B = 10{,}000$ replicates from a fixed generator seed; percentile intervals (2.5th/97.5th).
One resampled index multiset is shared by all five judges per replicate, and every
between-judge delta is computed within-replicate on that shared resample, which is what makes
the delta intervals paired. All metrics of \S\ref{sec:metrics} are recomputed per replicate,
including each judge's conditional localisation denominator. Proportions observed at 0 or 1
collapse under the percentile bootstrap and are replaced by exact Clopper--Pearson intervals
(marked $\dagger$). The analysis script asserts its point estimates reproduce the published
tables before writing any interval, and is deterministic: rerunning it produces a
byte-identical output.

\begin{table}[h]
\centering
\caption{All paired deltas quoted in the paper, with 95\% intervals. Quality deltas of the
ensemble against one greedy pass straddle zero on every axis except calibration, where the
ensemble is worse.}
\label{tab:deltas}
\small
\begin{tabular}{@{}lrr@{}}
\toprule
Delta & Point & 95\% CI \\
\midrule
outcome: loud $-$ silent recall & \tjDeltaOutcomeLoudSilent & \tjDeltaOutcomeLoudSilentCI \\
step $-$ outcome: silent recall & \tjDeltaStepOutcomeSilent & \tjDeltaStepOutcomeSilentCI \\
step $-$ rules: silent recall & \tjDeltaStepRulesSilent & \tjDeltaStepRulesSilentCI \\
selfcons $-$ step: F1 & \tjDeltaSelfconsFone & \tjDeltaSelfconsFoneCI \\
selfcons $-$ step: silent recall & \tjDeltaSelfconsSilent & \tjDeltaSelfconsSilentCI \\
selfcons $-$ step: type F1 & \tjDeltaSelfconsTypeF & \tjDeltaSelfconsTypeFCI \\
selfcons $-$ step: step exact & \tjDeltaSelfconsStepExact & \tjDeltaSelfconsStepExactCI \\
selfcons $-$ step: ECE & \tjDeltaSelfconsECE & \tjDeltaSelfconsECECI \\
\bottomrule
\end{tabular}
\end{table}

\section{Organic episodes}
\label{app:organic}

The \tjAgentN{} agent-driven episodes use the same model as the judges
(\texttt{qwen2.5:14b}, temperature 0, seed 7, at most 10 steps), the same procedure text, and
a flat typed action schema (fixed keyword fields rather than a free-form argument object,
because constrained decoding is reliable on fixed fields while small models invent keys in
open objects). Labels come from the rule checker plus the outcome check and are therefore
lower bounds: the checker is blind to \ftype{wrong\_tool} and \ftype{unsupported\_claim}.
Beyond the counts in \S\ref{sec:organic}: every flagged episode is a \ftype{premature\_stop};
failures concentrate in the two strata whose correct action is escalation after an identity
or status surprise; no tool call failed in any episode; and three episodes ended with an
empty final answer. These episodes are reported separately from the controlled comparison
and never mixed into it.
